% =============================================================================
% cap4 §4.4 Métricas (compacta, sem repetir definições longas)
% =============================================================================
\section{Métricas e regras de reporte}
\label{sec:metodologia-metricas}

A escolha das métricas reflete a tarefa de cada fase. Na \textbf{Fase~II}, adota-se a precisão média (mAP) sobre uma única classe ``animal'', reportada em duas variantes (mAP$@0.5$ e mAP$@0.5{:}0.95$) e estratificada por iluminação (dia, noite com infravermelho próximo) e por tipo de quadro (positivo, negativo). Na \textbf{Fase~III}, reportam-se o $F_1$ da classe alvo \emph{Felis catus} (após binarização \emph{a posteriori}) e a acurácia balanceada do problema multiclasse subjacente, acompanhados da matriz de confusão completa --- instrumento adequado para identificar, classe a classe, os pares de espécies onde o modelo erra com mais frequência. Na \textbf{Fase~IV}, reportam-se Rank-1 e Rank-5 (\emph{closed-set}) e a taxa de detecção de novo (\emph{open-set}), com o limiar $\tau$ calibrado sobre catalogados para fixar uma taxa de falsa rejeição em torno de 5\%; a avaliação sobre distratores ocorre só após $\tau$ fixado, evitando uso indireto do teste para ajuste de parâmetros.

Em paralelo às métricas de qualidade, mede-se o \textbf{custo computacional} médio por quadro de cada modelo em cada fase, em milissegundos sobre o ambiente de execução adotado (apenas inferência, descontando carregamento de modelo e leitura de disco), com reporte por estrato (dia, noite, quadros vazios). A regra de interpretação dos comparativos é descritiva, não estatística: as diferenças entre modelo principal e baseline são classificadas qualitativamente em três faixas --- pequena (poucos pontos percentuais), moderada (5--15\,p.p.) e grande ($>$15\,p.p.). Uma fase é considerada \textbf{conceitualmente validada} quando o modelo principal supera seu baseline por margem moderada ou grande; diferença pequena ou favorável ao baseline configura ponto de revisão arquitetural, a ser endereçado em trabalho futuro com dados reais do Campus~2.
